% ==================================================================
% id:       seq-002
% topic:    סדרות
% subtopic: הגדרת הגבול (ε, N)
% solution: full
% ==================================================================
\bE
\item
הוכיחו באינדוקציה כי לכל
$n$
טבעי מתקיים
$2^n\ge n$.
\hint{בצעד האינדוקציה כתבו $2^{n+1}=2\cdot2^n$ והשתמשו בהנחת האינדוקציה.}
\begin{solbox}
    נראה זאת באינדוקציה:
    \bit
\item
עבור
$n=1$,
מתקיים
$2^1\ge1$.
\item
נניח כעת כי
$2^n\ge n$,
ונראה כי
$2^{n+1}\ge n+1$.
ואכן מתקיים
$$.2^{n+1}=2\cdot 2^n
\mathrel{\mathop{\ge}\limits_{\makebox[0pt][c]{\scriptsize\text{האינדוקציה הנחת }}}}2\cdot n=n+n\ge n+1$$
    \eit
\end{solbox}

\item
בעזרת הגדרת הגבול הראו כי הסדרה
$a_n=\frac{1}{2^n}$
מתכנסת ל-$0$.

(בשלב מסוים היעזרו במה שהוכח בסעיף הקודם)
\hint{השתמשו בסעיף הקודם כדי לחסום את $\frac{1}{2^n}$ מלמעלה בביטוי פשוט יותר.}
\begin{solbox}
    יהי
    $\varepsilon>0$.
    נחפש
    $N$
    טבעי עבורו לכל
    $n>N$
    מתקיים
    $|a_n-0|<\varepsilon$.
    נסתכל על
    ההפרש האחרון, ונשתמש בסעיף (א):
    $$|a_n-0|=\left| \frac{1}{2^n}\right|\le
    \frac{1}{n}<\frac{1}{N}$$
לכן נבחר
$N$
המקיים
$\frac{1}{N}<\varepsilon$,
למשל נקח את
$$,N=\lfloor \frac{1}{\varepsilon}\rfloor +1$$
כך הראנו שלכל
$n>N$
מתקיים
$|a_n-0|<\varepsilon$,
כרצוננו.

\end{solbox}
\eE
